\section{Suite horizontal integral e integración vertical}

La sección anterior trató la diferenciación; esta trata la convergencia. En el mercado de consumo hay una tendencia clara a que todo converja en unos pocos líderes, y ChatGPT podría absorber muchos productos. El ejemplo de suite horizontal integral es ChatGPT, que reúne Chat, búsqueda, Coding, Agent y Workspace en una misma entrada; el ejemplo de integración vertical es Gemini, que, desde los chips TPU, el modelo Gemini y las aplicaciones Agent hasta Google Docs, Chrome, Android y YouTube, puede formar una superintegración.

\begin{figure}[H]
\centering
\includegraphics[width=0.94\textwidth]{horizontal-vs-vertical.png}
\caption{Suite horizontal integral frente a integración vertical: ChatGPT como suite horizontal, Gemini como superintegración vertical. Diagrama conceptual propio, elaborado a partir de la conversación de 00:21:37--00:33:35.}
\end{figure}

\begin{knowledgebox}{Lectura de la figura: lo horizontal y lo vertical son dos formas de absorber}
La suite horizontal integral concentra muchas tareas del usuario en una sola entrada; la integración vertical encadena chips, modelos, aplicaciones, navegador, sistema operativo y ecosistema de contenidos. La primera se apropia de la mente del usuario; la segunda, de la infraestructura y de las entradas predeterminadas.
\end{knowledgebox}

\subsection{La barrera de producto horizontal de ChatGPT}

ChatGPT no es solo un envoltorio de un modelo: tiene hábitos de los usuarios, posicionamiento de marca, una entrada predeterminada, historial de conversaciones y una expansión continua de funciones. Para un producto de IA independiente, lo más difícil es esto: si la función que construiste se convierte en una pestaña más de la suite de ChatGPT, ¿por qué el usuario seguiría usándote por separado?

\begin{importantbox}{El poder destructivo de la suite horizontal integral}
Lo más temible de la suite horizontal integral no es que cada función sea la mejor, sino que se convierte en la entrada predeterminada del usuario. Una entrada predeterminada puede distribuir con rapidez funciones «suficientemente buenas» a una enorme cantidad de usuarios y así acortar la ventana de oportunidad de los productos de las empresas emergentes.
\end{importantbox}

\subsection{La integración vertical de Gemini}

La ventaja de Google está en la integración vertical: TPU, nube, Gemini, búsqueda, publicidad, Docs, Chrome, Android y YouTube pueden funcionar de forma coordinada. Por eso mismo cambió la opinión de Guangmi (广密) sobre Google: cuando los productos de IA terminen por requerir la coordinación de búsqueda, publicidad, ofimática, navegador y sistema operativo, los activos antiguos de Google podrían volver a convertirse en ventajas nuevas.

\begin{knowledgebox}{Activos de la integración vertical de Google}
\begin{tabularx}{\textwidth}{p{0.22\textwidth}p{0.34\textwidth}X}
\toprule
Activo & Función & Valor en la era de la IA \\
\midrule
TPU/Cloud & Infraestructura de cómputo y de entrenamiento & Reduce el costo de los modelos y garantiza el suministro. \\
Search/Ads & Sistema de entrada y de monetización & La búsqueda con IA y las plataformas publicitarias podrían volver a fusionarse. \\
Chrome/Android & Navegador y sistema operativo móvil & Los Agent pueden entrar en el entorno digital predeterminado. \\
Docs/Workspace & Escenarios de oficina y archivos del usuario & Entrada a la productividad empresarial y personal. \\
YouTube & Contenido de video y datos/escenarios multimodales & Campo de aplicación para la comprensión y la generación multimodales. \\
\bottomrule
\end{tabularx}
\end{knowledgebox}

\subsection{Resumen de la sección}

La suite horizontal integral y la integración vertical explican por qué cada vez es más difícil esquivar a las empresas líderes. Los emprendedores deben convertirse en una capacidad clave dentro de la suite o encontrar escenarios que a los gigantes no les resulte fácil integrar.
